\documentclass{beamer}
\usetheme{AnnArbor}

\usepackage[T1]{fontenc}
\usepackage[utf8]{inputenc}
\usepackage[spanish]{babel}
\usepackage{amsmath}

\title{FPS, \(dt\)}
\subtitle{Animación: Movimiento discreto en Raylib}
\author{HCI}
\date{April 2026}

\begin{document}

\begin{frame}
\titlepage
\end{frame}

\begin{frame}{Objetivo}
Comprender con claridad:

\begin{itemize}
    \item Qué son los FPS.
    \item Qué es \(dt\).
    \item Qué ocurre al mover objetos sin \(dt\).
    \item Qué ocurre al mover objetos con \(dt\).
\end{itemize}
\end{frame}

\begin{frame}{Contexto del problema}
Queremos mover un círculo sobre el eje \(x\).

\medskip

La posición horizontal del círculo será:

\[
x(t)
\]

donde:

\begin{itemize}
    \item \(x(t)\): posición horizontal en el tiempo \(t\).
    \item \(t\): tiempo medido en segundos.
    \item \(v\): velocidad horizontal.
\end{itemize}
\end{frame}

\begin{frame}{¿Qué son los FPS?}
FPS significa:

\[
\text{Frames Per Second}
\]

Es decir:

\[
\text{Cuadros por segundo}
\]

\medskip

Si un programa corre a 60 FPS, intenta dibujar:

\[
60 \text{ imágenes cada segundo}
\]

En animación, cada imagen se llama \textbf{frame}.
\end{frame}

\begin{frame}{FPS y tiempo entre frames}
Relación ideal:

\[
dt \approx \frac{1}{\text{FPS}}
\]

\medskip

Ejemplos:

\[
30\text{ FPS} \Rightarrow dt \approx 0.0333
\]

\[
60\text{ FPS} \Rightarrow dt \approx 0.0167
\]

\[
120\text{ FPS} \Rightarrow dt \approx 0.0083
\]

A más FPS, menor duración de cada frame.
\end{frame}

\begin{frame}{¿Qué es \(dt\)?}
En simulación, \(dt\) significa:

\[
dt = \Delta t
\]

Es el tiempo transcurrido entre dos frames consecutivos.

\medskip

En Raylib se obtiene con:

\[
\texttt{GetFrameTime()}
\]

\medskip

Si el frame tarda mucho, \(dt\) aumenta.  
Si el frame tarda poco, \(dt\) disminuye.
\end{frame}

\begin{frame}{\(dt\) como puente con el tiempo real}
El computador actualiza la escena cuadro a cuadro.

\medskip

Pero el movimiento físico depende del tiempo real.

\medskip

Por eso usamos:

\[
\Delta x = v \cdot dt
\]

donde:

\begin{itemize}
    \item \(\Delta x\): desplazamiento en un frame.
    \item \(v\): velocidad en píxeles por segundo.
    \item \(dt\): segundos transcurridos en ese frame.
\end{itemize}
\end{frame}

\begin{frame}{Modelo matemático del movimiento}
El modelo continuo de velocidad constante es:

\[
\frac{dx}{dt} = v
\]

Esto significa:

\begin{itemize}
    \item La posición cambia con el tiempo.
    \item La velocidad \(v\) mide qué tan rápido cambia \(x\).
\end{itemize}

Si \(v=200\), el objeto debe avanzar:

\[
200 \text{ píxeles en un segundo}
\]
\end{frame}

\begin{frame}{Actualización correcta}
Para simular el movimiento, usamos:

\[
x_{\text{nuevo}} =
x_{\text{actual}} + v \cdot dt
\]

Esto significa:

\begin{itemize}
    \item Tomamos la posición actual.
    \item Calculamos cuánto debe avanzar.
    \item Sumamos ese avance a la posición.
\end{itemize}

Esta es una forma básica del método de Euler.
\end{frame}

\begin{frame}{Movimiento sin \(dt\)}
El movimiento incorrecto usa una constante fija:

\[
x_{\text{nuevo}} =
x_{\text{actual}} + v \cdot c
\]

En el código:

\[
c = 0.016
\]

Esto supone artificialmente que todos los frames duran:

\[
0.016 \text{ segundos}
\]

Esto solo se parece a 60 FPS.
\end{frame}

\begin{frame}{Problema de no usar \(dt\)}
Si el programa corre a 30 FPS, hay menos actualizaciones por segundo.

\medskip

Si corre a 120 FPS, hay más actualizaciones por segundo.

\medskip

Usando una constante fija:

\[
\text{distancia por segundo}
=
\text{actualizaciones} \times v \cdot c
\]

Por eso el movimiento depende de los FPS.
\end{frame}

\begin{frame}{Ejemplo sin \(dt\)}
Supongamos:

\[
v = 200,\qquad c=0.016
\]

Avance por frame:

\[
\Delta x = 200(0.016)=3.2
\]

A 30 FPS:

\[
30(3.2)=96
\]

A 120 FPS:

\[
120(3.2)=384
\]

El objeto no se mueve igual.
\end{frame}

\begin{frame}{Movimiento con \(dt\)}
Con \(dt\), cada frame usa su duración real:

\[
x_{\text{nuevo}} =
x_{\text{actual}} + v \cdot dt
\]

A 30 FPS:

\[
30 \cdot 200 \cdot 0.0333 \approx 200
\]

A 120 FPS:

\[
120 \cdot 200 \cdot 0.0083 \approx 200
\]

El movimiento respeta el tiempo real.
\end{frame}

\begin{frame}{Comparación conceptual}
\begin{center}
\begin{tabular}{|c|c|c|}
\hline
Caso & Fórmula & Consecuencia \\
\hline
Sin \(dt\) & \(x=x+v\cdot c\) & Depende de FPS \\
\hline
Con \(dt\) & \(x=x+v\cdot dt\) & Estable \\
\hline
\end{tabular}
\end{center}

\medskip

\[
dt \text{ corrige la duración real de cada frame.}
\]
\end{frame}

\begin{frame}{Integración numérica: idea básica}
El modelo continuo es:

\[
\frac{dx}{dt}=v
\]

Pero el computador no calcula todos los tiempos.

\medskip

Calcula una secuencia discreta:

\[
t_0,\; t_1,\; t_2,\; \ldots,\; t_n
\]

La separación entre dos tiempos es:

\[
\Delta t = t_{n+1}-t_n
\]
\end{frame}


\begin{frame}{Lectura del programa}
El programa compara dos modelos:

\begin{itemize}
    \item Círculo rojo: usa una constante fija.
    \item Círculo azul: usa el \(dt\) real.
\end{itemize}

Al cambiar FPS:

\begin{itemize}
    \item El rojo cambia su rapidez visible.
    \item El azul conserva su rapidez.
\end{itemize}
\end{frame}

\begin{frame}{Conclusión}
Ideas principales:

\begin{itemize}
    \item FPS indica cuántos frames se dibujan por segundo.
    \item \(dt\) mide cuánto duró el último frame.
    \item Sin \(dt\), el movimiento depende de los FPS.
    \item Con \(dt\), el movimiento respeta el tiempo real.
    \item La fórmula \(x=x+vdt\) es integración numérica.
    \item La integral continua produce \(x(t)=x_0+vt\).
\end{itemize}
\end{frame}

\end{document}
